% !TEX program = xelatex
\documentclass[11pt,a4paper]{article}

\usepackage[a4paper,margin=22mm]{geometry}
\usepackage{xeCJK}
\usepackage{fontspec}
\usepackage{amsmath,amssymb}
\usepackage{booktabs}
\usepackage{tabularx}
\usepackage{array}
\usepackage{enumitem}
\usepackage{xcolor}
\usepackage{hyperref}

\setCJKmainfont{Hiragino Mincho ProN}
\setCJKsansfont{Hiragino Sans}
\setCJKmonofont{Hiragino Sans}
\setmonofont{Menlo}[Scale=0.82]

\xeCJKDeclareCharClass{CJK}{"2460 -> "24FF, "2013 -> "2015, "2018 -> "201F, "2026, "2500}

\hypersetup{colorlinks=true, linkcolor=black, urlcolor=blue!60!black,
  pdftitle={6軸アームの制御に Xinu メッシュは効くか},
  pdfauthor={児玉靖司}}

\newcommand{\code}[1]{\texttt{#1}}
\newcommand{\prob}[1]{\textbf{\textcolor{red!55!black}{#1}}}
\newcommand{\good}[1]{\textbf{\textcolor{green!45!black}{#1}}}

\title{6軸アームの制御に Xinu メッシュは効くか\\
       \large --- 現状評価、効果の見積もり、実機構成の提案}
\author{}
\date{2026年7月21日}

\begin{document}
\maketitle

\begin{abstract}
\noindent
6 関節のアーム型ロボットを Raspberry Pi 5 経由で制御する構成において、
Xinu ベースの 12 台メッシュネットワークが効果を発揮するかを検討する。
結論を先に述べる。\prob{関節サーボループに対しては効果がなく、むしろ有害である}。
一方で\good{動作計画・学習・分散知覚・冗長性には効く}が、効果は台数に比例しない。
既存の実測値からの外挿では 12 ノードで 4--6 倍程度である。
さらに、\prob{現状の Xinu は 3 ボードとも I2C・PWM・GPIO ドライバを一切持たず、
アームに物理的に接続されていない}。本書はその根拠を示し、
メッシュの価値を最大化する実機構成を提案する。
\end{abstract}

\section{前提と問い}

対象は 6 関節のアーム型ロボットで、全ての関節指令は Pi 5 を通して送る構成を想定する。
本書が答える問いは 3 つ。

\begin{enumerate}[leftmargin=1.6em]
\item この構成で Xinu メッシュ 12 台は効果を発揮するか。
\item 効果があるとして、どの程度か。
\item メッシュの効果を最大限に生かすなら、どのような実機構成が考えられるか。
\end{enumerate}

評価対象の現物は Raspberry Pi 3 B+ / Pi 4 / Pi 5 の 3 台で、
WiFi ad-hoc (IBSS) による MANET を構成している
(SSID \code{MANET}、ch 6、静的 \code{10.0.0.<n>}、HELLO は UDP:5000 の 2 秒周期)。
2026-07-21 時点で 3 ノード完全メッシュが成立している。

\section{結論}

\begin{center}
\begin{tabularx}{\textwidth}{@{}l l X@{}}
\toprule
\textbf{用途} & \textbf{効果} & \textbf{理由}\\
\midrule
関節サーボループ & \prob{効果なし} & 制御ループが Pi 5 内で閉じるため、他ノードはクリティカルパス外。
                                     加えて WiFi の共有媒体が Pi 5 のジッタを増やす\\
動作計画 (RRT* 等) & \good{効く} & ほぼ完全並列。ただし線形には伸びない\\
強化学習のロールアウト & \good{効く} & エピソードが独立。既存 TinyML RL 資産と相性が良い\\
分散知覚 & \good{効く} & 各ノードで特徴抽出し、特徴だけを送る\\
冗長性・フェイルオーバ & \good{効く} & 安全性の観点では最も価値が高い\\
\bottomrule
\end{tabularx}
\end{center}

\section{なぜ制御ループには効かないのか}

\subsection{構成上の必然}

制御ループが Pi 5 $\rightarrow$ サーボである限り、他の 11 台は
\textbf{クリティカルパスの外}にしか存在できない。ネットワークを渡らない処理に
ノードを足しても、そのループは 1 マイクロ秒も速くならない。
これは実装の問題ではなく構成の問題である。

\subsection{数値で見た不一致}

6 軸アームの関節ループは典型的に 100--1000\,Hz、周期 1--10\,ms で回す。
ジッタ許容はその 1 割程度、すなわち 0.1--1\,ms 程度である。

\begin{center}
\begin{tabular}{@{}l r l@{}}
\toprule
\textbf{測定対象} & \textbf{実測値} & \textbf{出典}\\
\midrule
Pi 4 のローカル RT スケジューラのジッタ & max 10\,$\mu$s & \code{/rtos-jitter} 10\,ms 周期・CPU hog 下、ミス 0/100\\
メッシュ経由の ICMP RTT (WiFi) & 3--27\,ms & Pi 5 $\rightarrow$ Pi 3 / Pi 4、2026-07-21\\
必要なジッタ許容 (1\,kHz ループ) & 0.1\,ms & 一般的な設計値\\
\bottomrule
\end{tabular}
\end{center}

ローカルのスケジューラは\good{十分すぎる性能}を持っている。max jitter 10\,$\mu$s は
1\,kHz ループの許容値の 1/10 であり、RTOS として本物である。
問題はネットワークで、\prob{WiFi 経由のレイテンシは要求より 1--2 桁大きい}。
しかも CSMA/CA は競合ベースで再送を伴うため分布の裾が重く、平均値が満たせても
最悪値が保証できない。関節ループを WiFi 越しに回す設計は成立しない。

\section{効果の見積もり}

台数効果は既存の実測値から外挿できる。分散 N-Queens のベンチマークでは、
\textbf{3 ボード $\times$ 4 コア = 12 コアで N=14 を 1458\,ms}、最速単体比
\textbf{1.87 倍}であった。3 倍のハードウェアで 1.87 倍、すなわち\textbf{効率 62\%}である。
しかもこの問題は通信をほとんど必要としない、並列化に最も有利な部類に属する。

12 ノードに広げた場合、WiFi の帯域と結果を集約するコーディネータがボトルネックになるため、
\textbf{現実的な期待値は 4--6 倍}、通信を伴う問題ではそれ以下と見るべきである。

具体的な意味は次のとおり。単体で 500\,ms かかる動作計画が \textbf{100\,ms 前後}になる。
再計画の応答性としては意味のある改善だが、サーボループとは別次元の話である。

\section{現状の欠落}

\subsection{アームに繋がっていない}

両ボードの \code{device/} ディレクトリの内容は次のとおりである。

\begin{center}
\code{genet\ \ gic\ \ mbox\ \ sd\ \ timer\ \ uart\ \ usb\ \ video\ \ wifi}
\end{center}

\prob{I2C も PWM も GPIO もサーボドライバも存在しない}。
すなわち現状このメッシュは\prob{アームに物理的に一切接続されていない}。
DOFBOT や mecharm の成果は Mac 上のシミュレータ (AIPL / Three.js) であり、
実機の Xinu 側にアクチュエータ I/O は無い。

\subsection{ROS 相当として欠けているもの}

重要度順に挙げる。

\begin{enumerate}[leftmargin=1.6em,itemsep=3pt]
\item \prob{アクチュエータ・センサ I/O}。上記。これが無い限り何も動かない。最優先。
\item \prob{時刻同期}。各ノードが自分の \code{CNTPCT} を見ているだけで、ノード間の共通時刻がない。
      加えて時間軸自体に既知の欠陥がある --- rpi3 ではアクターランタイムが
      「100\,Hz クロック」を仮定して \code{clkticks * 10} しているが実際は 1000\,Hz で
      \textbf{10 倍ずれ}、さらに \code{clkticks} が毎秒 0 にリセットされ\textbf{単調でない}。
      センサ融合も協調軌道も成立しない。
\item \prob{Pub/Sub と QoS}。現在の通信は HTTP のリクエスト/レスポンスとアクターメッセージのみ。
      トピックによる多対多、信頼性/ベストエフォートの選択、デッドラインや liveliness がない。
\item \prob{有界遅延のトランスポート}。今の TCP は RTO 500\,ms の go-back-N で、制御ループには使えない。
      RT 制御には UDP ベースで再送回数を制限した別経路が要る。
\item \prob{座標系と運動学}。TF (フレーム木)、URDF 相当のロボット記述、FK/IK、
      軌道表現がボード側にない。
\item \prob{安全機構}。ウォッチドッグと非常停止がなく、ユーザ空間が存在しない
      (EL1 単一アドレス空間) ため、JIT で送り込んだコードの暴走がノード全体を落とす。
      アームを動かす前に必須。
\end{enumerate}

\subsection{あるもの}

公平のため、既に本物である部分も記す。Pi 4 のプリエンプティブ優先度スケジューラ
(実測 max jitter 10\,$\mu$s、100 回中ミス 0) は制御ループの土台として通用する。
マルチホップ MANET、アクターランタイム、ノードへコードを送り込む JIT も揃っている。
\textbf{土台はあるが、ロボットミドルウェアとしての層が丸ごと無い}、
というのが正確な評価である。

\section{実機構成の提案}

現在の「Pi 5 が 6 関節を直接叩く」構成は、\prob{メッシュの価値を最も引き出せない形}である。
メッシュを活かすなら、産業用アームが実際に採っている構造に寄せるのが筋である。

\subsection{提案 1: スマートジョイント化 (最重要)}

各関節にモータドライバと同居する小さな計算ノードを置き、
\textbf{電流・位置ループを関節ローカルで 1--10\,kHz} で閉じる。
ネットワークを渡るのは目標値と軌道だけ (100\,Hz--1\,kHz) とする。
こうして初めて「6 関節ノード + 計画ノード + 知覚ノード」という構成に意味が生まれる。

関節側は Pi である必要はない。\textbf{ハードウェア PWM とエンコーダ入力を持つ
RP2040 / Pico 2 や STM32} が適任である (RP2040 なら PIO でエンコーダを読める)。
これは EtherCAT スレーブを関節ごとに置く産業用アームと同じ構造である。

\subsection{提案 2: 決定性のあるバックボーン}

関節チェーンは \textbf{CAN-FD または EtherCAT}、上位は
\textbf{スイッチド Gigabit Ethernet}、可能なら \textbf{TSN (802.1Qbv)} とする。
WiFi ad-hoc は制御には向かない。

\subsection{提案 3: 無線メッシュが正当化される構成}

WiFi メッシュが本当に必要なのは、\textbf{物理的に離れたノードがある場合だけ}である。

\begin{itemize}[leftmargin=1.6em,itemsep=2pt]
\item \textbf{移動マニピュレータ} (台車 + アーム) --- 台車と外部センサ間は無線が必然
\item \textbf{複数アームの協調} --- 別筐体間の通信
\item \textbf{環境設置型センサ} (部屋のカメラ、触覚マット) --- 配線したくない場所
\end{itemize}

すなわち\textbf{アーム内部は有線・決定性、アームと外界は無線メッシュ}という役割分担が、
今の資産を最も活かす形である。

\section{今の機材で効果を測る実験}

新しいハードウェアを購入する前に、既存の 3 ボードで答えを出せる。

\begin{enumerate}[leftmargin=1.6em,itemsep=3pt]
\item 6 関節は予定どおり Pi 5 に直結する (まず I/O ドライバが必要)。
\item \textbf{動作計画だけをメッシュに出す}。Pi 5 が計画クエリを投げ、
      $N$ ノードが並列に探索し、最良の軌道を返す。
\item 単体ノード時とのエンドツーエンド計画レイテンシを比較する。
\end{enumerate}

これは既存の JIT とアクターランタイムで\textbf{新規ハードウェアなしに測定できる}。
N-Queens で 1.87 倍という実測値が既にあるので、同じ枠組みで計画問題の実効値が取れる。
この数字が出て初めて、12 台に増やす価値があるかを根拠を持って判断できる。

\section{本書の測定値について}

数値の出所と限界を明示する。

\begin{itemize}[leftmargin=1.6em,itemsep=2pt]
\item \textbf{max jitter 10\,$\mu$s / ミス 0/100} --- Pi 4 実機、\code{/rtos-jitter}、
      10\,ms 周期・CPU hog 下での実測。
\item \textbf{メッシュ RTT 3--27\,ms} --- 2026-07-21 に Pi 5 から Pi 3 / Pi 4 へ実施した
      ICMP の実測。サンプル数は少なく、\textbf{最悪値の統計は取っていない}。
      制御用途の判断には本来もっと長時間の分布測定が要る。
\item \textbf{N-Queens 3 ボード 1.87 倍} --- 既存ベンチマークの記録値。
      \textbf{12 ノードでの 4--6 倍は実測ではなく外挿}であり、
      WiFi の帯域とコーディネータの負荷という 2 つの仮定に依存する。
\item \code{device/} の内容は 2026-07-21 時点のソースツリーを直接確認した結果である。
\end{itemize}

\vfill
{\small\itshape 対象: xinu-rpi3 (BCM2837 / Cortex-A53)、xinu-rpi4 (BCM2711 / Cortex-A72)、
xinu-rpi5 (BCM2712 / Cortex-A76)。メッシュは WiFi ad-hoc (IBSS)、SSID \code{MANET}、
ch 6、\code{10.0.0.2/3/5}。}

\end{document}
